% Portable starter following the public Revista Brasileira de Economia layout.
\documentclass[11pt,a4paper]{article}
\usepackage[utf8]{inputenc}
\usepackage[T1]{fontenc}
\usepackage[brazil]{babel}
\usepackage[margin=2.6cm]{geometry}
\usepackage{amsmath,booktabs,tabularx,xcolor,hyperref,microtype}
\hypersetup{colorlinks=true,linkcolor=blue!50!black,urlcolor=blue!50!black}
\title{Produtividade e adoção gradual de ferramentas digitais}
\author{Camila Souza\thanks{Departamento de Economia, Universidade de Exemplo}\and
  Bruno Almeida}
\date{}

\begin{document}
\maketitle
\begin{center}
\small\textbf{Códigos JEL:} D24, O33, C23\quad
\textbf{Palavras-chave:} produtividade; tecnologia; firmas
\end{center}

\begin{abstract}
Este artigo apresenta um exemplo de como documentar uma análise empírica sobre
a adoção de ferramentas digitais por firmas. A especificação separa o efeito
de produtividade da tendência temporal e explicita as limitações dos dados.
Os resultados numéricos são fictícios e servem apenas para demonstrar a
estrutura do manuscrito.
\end{abstract}

\section{Introdução}
A adoção de uma nova ferramenta pode reduzir custos, mas também exige tempo de
aprendizagem. Uma análise econômica deve, portanto, registrar o período de
implementação e as características das firmas comparadas.

\section{Estratégia empírica}
Considere a especificação de diferenças em diferenças
\begin{equation}
 y_{it}=\alpha_i+\gamma_t+\beta(Trat_i\times Pos_t)+\varepsilon_{it},
 \label{eq:did}
\end{equation}
na qual $\beta$ mede a variação média associada à adoção após o período de
implementação. Em uma aplicação real, a identificação exige uma discussão
explícita da tendência paralela e das decisões de amostragem.

\section{Resultados}
\begin{table}[ht]
\centering
\caption{Efeito ilustrativo da adoção digital.}
\label{tab:effect}
\begin{tabular}{lrr}
\toprule
Especificação & Coeficiente & Erro-padrão\\
\midrule
Sem controles & 0.084 & 0.031\\
Controles básicos & 0.067 & 0.029\\
Efeitos fixos & 0.052 & 0.027\\
\bottomrule
\end{tabular}
\end{table}
Os valores da Tabela~\ref{tab:effect} não constituem evidência empírica; eles
mostram como apresentar estimativas e incertezas de forma compacta.

\section{Discussão}
O resultado pode refletir seleção das firmas, mudanças de composição ou
choques simultâneos. Um manuscrito completo deve incluir a fonte dos dados,
limpeza, código de estimação e testes de robustez.

\section{Conclusão}
Uma apresentação transparente facilita a avaliação da hipótese econômica e
permite que o leitor substitua o exemplo por seus próprios dados. Mantenha os
metadados do artigo e as notas de replicação junto ao arquivo final.

\end{document}
